\section{遮光カーテンについて/already}
\rightline{（文責：千のタリスマン）}

遮光カーテンについて

\subsection{1}
june気づいてくれなかったこと。傘に萌芽するあなたは、誕生するあなたはすべて関係している。廊下の長さが少し伸びる。
\begin{center}
$\ast \cdot \dots \cdot \star \cdot \dots \cdot \ast$
\end{center}

\subsection{2}
かなりの量の弱点が存在していて、各々がそれらを隠しながらビリヤニやハッピー・ターンなどを食べて生活を続かせている。最近はカルディでも弱点ガードなるプラスチックのお椀のようなものを売り出しているっぽい。質を求めるのならロフトの方が良いとか……

\subsection{3}
ロケットは発射しているのではなく、ゆっくり進化をしているのと同じで本は開かれているのではなく失くされている。パニーニは食べるものではなく打破するものであるという大胆な発想から話題になった、徳光和夫著『黒いお皿に置いてパニーニは打破しなさい』（大竹書房、現在絶版）がネットオークションで高値で取引されているニュースは皆の記憶に新しいだろう。とにかく、そういうことなのだ。

\subsection{4}
書面でなされた光に関する約束は意味をなさない $\spadesuit$ \\
あらゆる数学の定理は白色のインクで書かれたとき破綻する $\heartsuit$ \\
ハイライズ・ビルディングスは正確には「窓」に分類される $\diamondsuit$

\begin{center}
--- $\dagger$ ---
\end{center}

\subsection{already}
押入れから手紙が届いた。もう限界だという。
よく晴れた冬の朝、珍しくすっぱり床から這い上がったところ、枕元に一通、手紙が差し込まれていた。差し込まれていた。落ちていたという表現は適切でなく、置かれていたというのとも異なり、挟まるほどの圧力はなく、投函するには場所がない。枕元に差し込むというのが絶妙で、枕にでもなく布団にでもなく枕元。そばと言うには範囲が狭い。ともかく枕元に差し込まれていた。そういうほかない。否定を並べて示すくらいが限界だ。
ともかく、手紙のことだ。宛先、僕の家。送り主、僕の家。もちろん切手はないし消印もない。可愛らしい便箋にシンプルな文面がひとつ。
「もう限界です」
押し入れは二段構造になっていて、襖が二つで部屋から仕切る。機械の同居人がいるわけでもなく、とくに居心地は良くない。乱暴に服が突っ込まれるばかりで、空間はなく、だから万年床があるのだが、それが迷惑かと問われて、そちらの踏み入る領分ではないだろうと思う。罪の意識もない。
「なにが」
ととりあえず問うてみる。がらんどうの内に虚しく響く。返答など戻るはずもなく、話せるならわざわざこんな回りくどいことをする必要はないとそれが道理だ。
襖を優しくフェザーに触る。丸く凹んだノブに手をかけ、右に動かす。ゆっくり加速し、木が軋み砂の挟まる音がやけに耳につく。
減速。静止するのと端に到達する瞬間が同時に起こり、これはちょっとした僕の特技だ。日常生活の一々を儀式めいて行うことでスターになりきる。朝日が丸く、僕に当たる。
「え」
もちろん何もない。埃はある。神妙な顔のまま咳き込んでみる。埃が舞う。視界が白い。ようやく晴れてその先には「あ」
もちろん何もない。腐るということもなく元気そうにしている。築十年。まだまだこれからだ。
飽きて閉める。
布を3枚整えて、間へ僕を差し込む。去年から電気ストーブは壊れたまま部屋の隅に転がっている。押し入れとは逆の側なので文句のつけようはないはずで、今は冬だ。
悪い夢なら覚めてみろと、これが悪い夢がどうかに自信はない。良い夢でないのは確かだが、怖さと比べてどちらかといえば興味が勝つ。勝ったところで一寸法師の背比べで、特に意欲は湧かない。こうして二度寝へ向かっているのがいい証拠だ。

\begin{center}
$\approx$
\end{center}

こんな夢を見る。
魚屋の前にいて、店主と頬を張り合う夢だ。互いに頬を張り合い、拳を突き合せ、抱擁を交わす。磯の匂いのする胸を突き飛ばして、僕らは微笑を携えて顔を見詰める。するとそのうち魚の焦げる匂いがして、僕は鏡を見ていたことに気付く。
あるいは八百屋の前で、トマトを投げつけ合う夢。人参を店主の鼻に突き刺して、僕はもやしを髪にまぶす。じゃがいもを放りあって当たった部位に応じて数字を叫ぶ。こうしてセロリを購入するのだ。

\begin{center}
$\approx$
\end{center}

ずいぶん長く夢を見ていた。
雨のしとしと落ちる冬の夜、いつものように、彼は死体だった。